\documentclass[5pt, a4paper]{article}

\usepackage[utf8]{inputenc}
\usepackage[T1]{fontenc}
\usepackage{textcomp}
\usepackage{amsmath, amssymb}
\usepackage{multicol}
\usepackage{proof}
\usepackage{enumitem}
\usepackage{fancyhdr}
\usepackage[margin=1in]{geometry}
% figure support
\usepackage{import}
\usepackage{xifthen}
\pdfminorversion=7

\usepackage[most]{tcolorbox}
\usepackage{pdfpages}
\usepackage{transparent}
\newcommand{\incfig}[1]{%
    \def\svgwidth{\columnwidth}
    \import{./figures/}{#1.pdf_tex}
}

\tcbset{
colback=gray!5,
colframe=blue!50!white,
coltitle=black,
arc=0pt,
boxrule=-0.1mm,
fonttitle=\bfseries,
coltitle=black,
left=2.5mm,
leftrule=1mm,
right=3.5mm,
colbacktitle=black!10!white,
toptitle=0.75mm,
bottomtitle=0.5mm,
}
\begin{document}
\begin{tcolorbox}[title=Question 1]
    Compare the cost of the following computations: \[
        (x x^{T})A \quad \text{and} \quad x(x^{T}A)
    .\] 
\end{tcolorbox}
For the first one, we perform a vector multiplication that would lead into a $n
\times n$ amtrix. This leaves us with $n^2$ operations. A square matrix
multiplication, using the naive algorithm, would yield $O(n^2(2n-1))$.

For the second one, we have $x^{T}A$, which yields $n(2n - 1)$. Computing the
next one woud yield $n^2$ operations. 

Thus, calculating the result using the order on the right would be
computaionally less expensive ($n^2$) by a lot.
\begin{tcolorbox}[title=Question 2]
    Consider the matrix \[
        D_n =
   \begin{pmatrix} 
       1 & -2 &  & &  \\
        & 1 & -2 & &  \\
        &  & \ddots & \ddots &  \\
        &  & & 1 & -2 \\
        &  & &  & 1 \\
   \end{pmatrix}  
    .\] 
    Write out the general formula for $D_i^{-1}$
\end{tcolorbox}
For $n = 2$
 \[
     D_2^{-1} = \begin{pmatrix} 1 & 2 \\ 0 & 1 \end{pmatrix} 
.\] 
For $n = 3$,
 \[
D_3^{-1} = 
\begin{pmatrix}  
    1 & 2 & 4 \\ 
    0 & 1 & 2 \\
    0 & 0 & 1 \\
\end{pmatrix} 
.\] 
We can generalize this with \[

.\] 

\begin{tcolorbox}[title=Question 3]
\end{tcolorbox}
\begin{tcolorbox}[title=Question 4]
\end{tcolorbox}
\begin{tcolorbox}[title=Question 5]
\end{tcolorbox}
\begin{tcolorbox}[title=Question 6]
\end{tcolorbox}
\begin{tcolorbox}[title=Question 7]
\end{tcolorbox}
\end{document}
